\section{Primary powers and their intrinsic torsion}
\label{sec:intrinsic-primary-powers}
Fix $R\in G_e^\circ$ and abbreviate $S=S_R$, $X=X_R$,
$U=U_R$, $\Delta=\Delta_R$. On $\Delta$ let
$H=\operatorname{im}(K\to V)$ and $\Lambda=K\cap S$; these are
bundles of ranks $e-1$ and one. Define the closed subscheme
\begin{equation}\label{eq:ramification-embedded-support}
 E=\Delta\times_{\Pj(V^*)}Y_R\subset U.
\end{equation}
All ideals in this section are ideals of the displayed ambient scheme,
not merely ideals of set-theoretic supports.

\begin{lemma}[A global two-stratum primary calculation]
\label{lem:global-two-stratum}
Let $U$ be a smooth complex variety, let $\Delta\subset U$ be a
smooth integral Cartier divisor, and let $E\subset\Delta$ be a
nonempty smooth integral Cartier divisor of $\Delta$. Put
$L=\mathcal I_\Delta$ and $P=\mathcal I_E$. For every $n\ge1$,
\begin{equation}\label{eq:global-two-stratum}
 L^nP^n=L^n\cap P^{2n}.
\end{equation}
This is an irredundant primary decomposition with associated points
exactly $\eta_\Delta,\eta_E$. The specified primary representatives
$L^n$ and $P^{2n}$ are intrinsic ordinary powers of the two ideals.
\end{lemma}
\begin{proof}
At a point of $E$, choose local equations $t,f$ with
$L=(t)$ and $P=(t,f)$. They are part of a regular system of
parameters. The associated graded ring of $P$ is locally a
polynomial ring in two variables over the domain $\mathcal O_E$.
Multiplication by an element outside $P$ is injective on each of
its graded pieces, so induction on the filtration proves that
$P^a$ is $P$-primary for every $a\ge1$. The initial form of $t$
is one of the two polynomial variables; multiplication by it is
injective on the graded ring. It follows, by comparing the first
nonzero graded piece, that
\[
 (P^{2n}:t^n)=P^n.
\]
Hence $(t^n)\cap P^{2n}=t^nP^n$. The same argument for the
principal smooth divisor proves that $(t^n)$ is $(t)$-primary.
Off $E$, the equality reduces to $L^n=L^n$. These local identities
glue because all displayed ideals are intrinsic.

At $\eta_E$, the element $t^n$ lies in $L^n$ but not in $P^{2n}$,
whereas $f^{2n}$ lies in $P^{2n}$ but not in $L^n$. Neither
component can be omitted. A noetherian irredundant primary
decomposition with these two distinct radicals has exactly these
associated primes. The calculation on affine opens gives the
associated points stated globally. In particular, this proof does
not replace an equality of ideals by a claim about radicals.
\end{proof}

\begin{theorem}[Ramification as an embedded primary support]
\label{thm:ramification-primary}
For every $R\in G_e^\circ$ and every $m\ge2$, the multiplication
failure scheme $D=D_R$ on all of $U$ has ideal
\begin{equation}\label{eq:ramification-primary-ideal}
 \mathcal I_D=\mathcal I_\Delta\mathcal I_E
             =\mathcal I_\Delta\cap\mathcal I_E^2.
\end{equation}
The two supports are smooth and integral, and are its complete set
of associated supports on $U$. There are natural identifications
\begin{align}
 \Delta&\simeq
 \operatorname{Tot}_{\Pj(V^*)\times\Pj(S)}
                     \operatorname{Hom}(H,S/\Lambda),
       \label{eq:delta-graph-bundle}\\
 E&\simeq
 \operatorname{Tot}_{Y_R\times\Pj(S)}
                     \operatorname{Hom}(H,S/\Lambda).
       \label{eq:embedded-graph-bundle}
\end{align}
The vector-bundle rank is $(e-1)(p-1)$; $\dim U=ep$,
$\dim\Delta=ep-1$, and $\dim E=ep-2$.
The nilradical is square-zero and fits into the exact sequence
\begin{equation}\label{eq:ramification-nil-sequence}
 0\longrightarrow\mathcal O_E(-\Delta)
 \longrightarrow\mathcal O_D
 \longrightarrow\mathcal O_\Delta\longrightarrow0.
\end{equation}
\end{theorem}
\begin{proof}
For $K\in\Delta$, its image is a hyperplane $H$ and its
intersection with $S$ is a line $\Lambda$. The subspace $K/\Lambda$
is the graph of a unique map $H\to S/\Lambda$. Conversely the
inverse image of any such graph in $H\oplus S$ is an $e$-plane
with the prescribed image and intersection. This construction and
its inverse commute with base change; it proves
\eqref{eq:delta-graph-bundle} as a scheme isomorphism, not just a
bijection of points. Restricting to $Y_R$ proves
\eqref{eq:embedded-graph-bundle}. Smoothness, integrality, and the
dimensions follow. The rank-$e-1$ determinantal divisor is smooth
on $U$, since an $(e-1)$-minor is invertible there.

We compute the \emph{full} multiplication ideal. Since the
augmentation ideal has cube zero, the image of every symmetric
power of degree $m\ge2$ is $\C\cdot1+K+K^2$. On a frame chart
write $K\to V\oplus S$ as $\binom{M}{A_0}$. Its multiplication
presentation has scalar, linear and quadratic blocks
\[
 \begin{pmatrix}
 1&0&0\\ 0&M&0\\0&A_0&\gamma_R\Sym^2M
 \end{pmatrix}.
\]
Every nonzero maximal minor uses the scalar column and all $e$
columns meeting the $V$-rows. Expansion gives
\[
 \mathcal I_D=(\det M)I_p(\gamma_R\Sym^2M).
\]
This argument is valid over the chart's coefficient ring; it is
not a fibrewise rank test.

Near $\Delta$, invert an $(e-1)$-minor and apply invertible row
and column operations to make $M=\diag(I_{e-1},t)$. The moving
basis of $V$ decomposes the quadratic matrix into
\[
 [A,\ tB,\ t^2c],\qquad A=\gamma_R|_{\Sym^2H}
              \text{ along }t=0.
\]
The condition $\gamma_R(HV)=S$ holds everywhere by
Proposition~\ref{prop:quadratic-ramification}. Locally $[A,B]$
is therefore onto, so there are matrices $U_0,V_0$ with
$AU_0+BV_0=I_p$. It follows that the column module of
$[A,tB,t^2c]$ equals that of $[A,tI_p]$: multiplying this identity
by $t$ gives one inclusion, and the reverse inclusion is immediate.
Zeroth Fitting ideals agree for presentations with the same image.

If $A$ is invertible, the ideal is $(t)$. If its restriction
lies on $Y_R$, it has rank $p-1$. Inverting a nonzero $(p-1)$-minor
and taking the Schur complement reduces the remaining quadratic
presentation to a single row with ideal $(f,t)$, where $f\bmod t$
is a local equation of $Y_R$. Thus
\begin{equation}\label{eq:ramification-local-model}
 \mathcal I_D=(t^2,tf)=t(t,f).
\end{equation}
The equations $t,f$ cut out $E$ and are part of a regular system
of parameters. Off $\Delta$, $M$ is invertible, $\gamma_R$ is
onto, and the failure ideal is the unit ideal. This proves
$\mathcal I_D=\mathcal I_\Delta\mathcal I_E$ everywhere on $U$.
Lemma~\ref{lem:global-two-stratum} gives the primary decomposition
and its irredundancy and associated points directly on $U$.
Alternatively the same identities descend by
Lemma~\ref{lem:frame-primary-descent}; the graph descriptions
specify all descended supports.

Finally the radical of $\mathcal I_D$ is $\mathcal I_\Delta$,
and this ideal is invertible. Its quotient by
$\mathcal I_\Delta\mathcal I_E$ is
$\mathcal I_\Delta\otimes\mathcal O_E=\mathcal O_E(-\Delta)$.
Its square vanishes because $\mathcal I_\Delta\subset\mathcal I_E$.
It is nonzero at $\eta_E$. This proves the exact sequence and
square-zero assertion.
\end{proof}

\begin{theorem}[All powers and canonical embedded torsion]
\label{thm:primary-powers-intrinsic}
Keep the hypotheses of Theorem~\ref{thm:ramification-primary}.
For $n\ge1$ put $D^{[n]}=V(\mathcal I_D^n)$ and
$E^{[n]}=V(\mathcal I_E^n)$, as subschemes of $U$. Then
\begin{equation}\label{eq:ramification-all-powers}
 \mathcal I_D^n=\mathcal I_\Delta^n\mathcal I_E^n
               =\mathcal I_\Delta^n\cap\mathcal I_E^{2n}
\end{equation}
is an irredundant primary decomposition. There are exactly two
associated points, $\eta_\Delta$ and $\eta_E$. The generic
multiplicity on $\Delta$ is $n$. The nilradical has exact index
$2n$ along $E$ and index $n$ on $\Delta\setminus E$.

Let $\mathcal T_n$ be the kernel of localization of
$\mathcal O_{D^{[n]}}$ at its unique minimal associated point.
This is an intrinsic coherent module, and
\begin{equation}\label{eq:canonical-embedded-torsion}
 \mathcal T_n\simeq\mathcal O_{E^{[n]}}(-n\Delta),\qquad
 V\!\left(\operatorname{Ann}_{\mathcal O_{D^{[n]}}}\mathcal T_n\right)
       =E^{[n]}.
\end{equation}
There is an exact sequence
\begin{equation}\label{eq:power-torsion-sequence}
 0\longrightarrow\mathcal O_{E^{[n]}}(-n\Delta)
 \longrightarrow\mathcal O_{D^{[n]}}
 \longrightarrow\mathcal O_{n\Delta}\longrightarrow0.
\end{equation}
At $\eta_E$ the module $\mathcal T_n$ has length $n(n+1)/2$.
For $n=1$ it is exactly the nilradical. Thus the embedded support,
its ordinary-power neighbourhoods, and these torsion lengths are
canonical, even though an embedded primary representative in an
arbitrary decomposition need not be unique.
\end{theorem}
\begin{proof}
The ideal product follows by taking ordinary powers; the intersection
and associated points follow from Lemma~\ref{lem:global-two-stratum}.
Near $\eta_\Delta$, $f$ is invertible, so the ideal is $(t^n)$
and its generic multiplicity is $n$. At a point of $E$, its radical
is $(t)$ and the ideal is $t^n(t,f)^n$. The element $t^{2n}$ lies
in this ideal, whereas $t^{2n-1}$ does not, since cancellation of
$t^n$ would give $t^{n-1}\in(t,f)^n$. The graded-ring argument in
the lemma excludes that membership. Off $E$ the nilradical index
is $n$; for $n=1$ the zero nilradical has index one by convention.

On an affine neighbourhood, localization at the minimal prime
$(t)$ makes $(t,f)$ the unit ideal. The contraction of the localized
failure ideal is $(t^n)$, which is $(t)$-primary. The localization
kernel is therefore
\[
 (t^n)/t^n(t,f)^n\simeq \mathcal O_U/(t,f)^n
\]
with the invertible factor $t^n$ recording the transition functions.
This proves the first formula in
\eqref{eq:canonical-embedded-torsion} and the exact sequence. Its
annihilator is $(t,f)^n/\mathcal I_D^n$, because multiplication
by $t^n$ can be cancelled in the regular ambient ring. Hence the
annihilator defines exactly $E^{[n]}$, with its full scheme structure.
The local calculations agree on overlaps, both by their ideal
formulas and by the intrinsic definition of the localization kernel.
At $\eta_E$ the ambient ring is regular local of dimension two,
with maximal ideal $(t,f)$. The monomials $t^if^j$, $i+j<n$,
give the successive graded pieces of its quotient by $(t,f)^n$,
whose total length is $1+2+\cdots+n$. The claim for $n=1$ follows
also from \eqref{eq:ramification-nil-sequence}.
\end{proof}

\begin{proposition}[Flat relative primary neighbourhoods]
\label{prop:ramification-flat-family}
Over $G_e^\circ$ the construction gives relative schemes
$\mathcal E\subset\boldsymbol\Delta\subset\mathcal U$ and
$\mathcal D\subset\mathcal U$. The morphism
$\mathcal E\to G_e^\circ$ is smooth. For every $n\ge1$,
$\mathcal E^{[n]}$ and $\mathcal D^{[n]}$ are flat over $G_e^\circ$,
and \eqref{eq:power-torsion-sequence} holds relatively and remains
exact after arbitrary base change. Each complex fibre has the
primary decomposition \eqref{eq:ramification-all-powers}.
\end{proposition}
\begin{proof}
Use the universal quotient bundle $\mathcal S$ on $G_e^\circ$;
no global trivialization of it is required. The universal smooth
ramification hypersurface is proper and smooth over $G_e^\circ$,
and the graph description makes $\mathcal E$ a vector bundle over
its product with $\Pj(\mathcal S)$. It is therefore smooth over
the base. The relative determinant divisor is smooth as well.
Locally along $\mathcal E$, the preceding elimination works over
the parameter ring and gives $\mathcal I_{\mathcal D}=t(t,f)$.
The equations $t,f$ are relative regular parameters: $t$ defines
the relative smooth divisor and $f$ its smooth relative
ramification hypersurface. In particular their conormal module
is locally free, and all symmetric powers of it are flat over
the parameter base.

The conormal filtration of $\mathcal O_{\mathcal E^{[n]}}$ has
successive quotients
$\Sym^j(\mathcal I_{\mathcal E}/\mathcal I_{\mathcal E}^2)$
for $0\le j<n$, viewed on $\mathcal E$. Each is base-flat; hence
$\mathcal O_{\mathcal E^{[n]}}$ is base-flat. The analogous
principal filtration proves flatness of
$\mathcal O_{n\boldsymbol\Delta}$. Cancellation of $t^n$ gives
the relative version of \eqref{eq:power-torsion-sequence}; its
outer terms are flat, so its middle term is flat. Tensoring this
sequence with any base algebra is exact because the last term is
flat. Relative regularity also identifies the base-changed ideals
with the specified products and powers. The primary statement
on every complex fibre is supplied by the direct fibre proof in
Lemma~\ref{lem:global-two-stratum}, not by an assertion that chosen
primary components specialize in arbitrary families.
\end{proof}
